% GENERATED FILE. Edit canonical lessons, then run npm run generate:books.
% canonical-source-hash: fnv1a64:d84ee5c653aafcb3
% canonical-lessons: MR-C475-sahsi, MR-C475-darpok, MR-C475-bindhasta, MR-C475-bephikir, MR-C475-bhine

\chapter{Daring and Fearful}
\label{ch:mr-daring-and-fearful}
\hlchaptermodality{475}

\begin{chapteropening}
\textbf{By the end of this chapter:} \emph{I can describe a daring, fearful or carefree person.}

Complete the last lesson of chapter 475: I can describe a daring, fearful or carefree person.
\end{chapteropening}

\section[\texorpdfstring{\mr{साहसी} (sāhsī)}{साहसी (sāhsī)}]{\mr{साहसी} (sāhsī) — adventurous}

\label{lesson:MR-C475-sahsi}



\begin{quote}
Before the new one: say the Marathi for a mess, then the Marathi for discipline.
\end{quote}

\subsection*{You'll want to know: \mr{साहसी}}
\textbf{\mr{साहसी}} — \emph{sāhsī} — \textquotedblleft{}adventurous\textquotedblright{}.

\begin{quote}
\mr{साहसी} \mr{मुलांनी} \mr{डोंगर} \mr{चढून} \mr{पाहिला}.
\end{quote}

\begin{quote}
\emph{sāhsī mulānnī ḍoṅgar caḍhūn pāhilā.}
\end{quote}

\begin{quote}
The adventurous boys tried climbing the mountain.
\end{quote}

From \mr{साहस}, a daring deed. Unchanging.

\begin{cousinweb}
\textbf{\mr{साहसी}} (\emph{sāhsī}) is a Marathi adjective built on \mr{साहस}, \textquotedblleft{}a daring deed\textquotedblright{}, a tatsama from Sanskrit \emph{sāhasa}, \textquotedblleft{}boldness\textquotedblright{} (CDIAL 13378).
\end{cousinweb}

\begin{grammarlens}[title={What you've built}]
One more word for describing what you lived and felt.
\end{grammarlens}

\subsection*{Your turn}
\begin{itemize}
\raggedright
  \item \emph{Say it:} \emph{sāhsī}
  \item \emph{Say it:} \emph{sāhsī}, once more
  \item \emph{Say it:} \emph{sāhsī mulānnī ḍoṅgar caḍhūn pāhilā.} — The adventurous boys tried climbing the mountain.
  \item \emph{Recall:} say the Marathi for a mess, then the Marathi for discipline, then say \emph{sāhsī} again
\end{itemize}

\begin{tcolorbox}[breakable,colback=teal!4,colframe=teal!35!black,title={Before you move on}]
What does \mr{साहसी} mean? (\textquotedblleft{}Adventurous\textquotedblright{}.) Say it once more.
\end{tcolorbox}
\section[\texorpdfstring{\mr{डरपोक} (ḍarpok)}{डरपोक (ḍarpok)}]{\mr{डरपोक} (ḍarpok) — chicken-hearted}

\label{lesson:MR-C475-darpok}



\begin{quote}
Before the new one: say the Marathi for discipline, then the Marathi for adventurous.
\end{quote}

\subsection*{You'll want to know: \mr{डरपोक}}
\textbf{\mr{डरपोक}} — \emph{ḍarpok} — \textquotedblleft{}chicken-hearted\textquotedblright{}.

\begin{quote}
\mr{मी} \mr{डरपोक} \mr{नाही}, \mr{पण} \mr{साप} \mr{पाहिला} \mr{की} \mr{घाबरतो}.
\end{quote}

\begin{quote}
\emph{mī ḍarpok nāhī, paṇ sāp pāhilā kī ghābarto.}
\end{quote}

\begin{quote}
I am no coward, but I get scared when I see a snake.
\end{quote}

Unchanging, and blunter than \mr{भित्रा}.

\begin{cousinweb}
\textbf{\mr{डरपोक}} (\emph{ḍarpok}) is a loan from Hindustani \emph{ḍarpok}, \textquotedblleft{}timid\textquotedblright{}, built on \emph{ḍar}, \textquotedblleft{}fear\textquotedblright{}.
\end{cousinweb}

\begin{grammarlens}[title={What you've built}]
One more word for describing what you lived and felt.
\end{grammarlens}

\subsection*{Your turn}
\begin{itemize}
\raggedright
  \item \emph{Say it:} \emph{ḍarpok}
  \item \emph{Say it:} \emph{ḍarpok}, once more
  \item \emph{Say it:} \emph{mī ḍarpok nāhī, paṇ sāp pāhilā kī ghābarto.} — I am no coward, but I get scared when I see a snake.
  \item \emph{Recall:} say the Marathi for discipline, then the Marathi for adventurous, then say \emph{ḍarpok} again
\end{itemize}

\begin{tcolorbox}[breakable,colback=teal!4,colframe=teal!35!black,title={Before you move on}]
What does \mr{डरपोक} mean? (\textquotedblleft{}Chicken-hearted\textquotedblright{}.) Say it once more.
\end{tcolorbox}
\section[\texorpdfstring{\mr{बिनधास्त} (bindhāsta)}{बिनधास्त (bindhāsta)}]{\mr{बिनधास्त} (bindhāsta) — carefree; unafraid}

\label{lesson:MR-C475-bindhasta}



\begin{quote}
Before the new one: say the Marathi for adventurous, then the Marathi for chicken-hearted.
\end{quote}

\subsection*{You'll want to know: \mr{बिनधास्त}}
\textbf{\mr{बिनधास्त}} — \emph{bindhāsta} — \textquotedblleft{}carefree; unafraid\textquotedblright{}.

\begin{quote}
\mr{परीक्षा} \mr{संपल्यावर} \mr{मुली} \mr{बिनधास्त} \mr{खेळल्या}.
\end{quote}

\begin{quote}
\emph{parīkṣā samplyāvar mulī bindhāsta kheḷlyā.}
\end{quote}

\begin{quote}
When the exams were over the girls played without a care.
\end{quote}

\mr{बिन}, without, + \mr{धास्त}, dread, the word in \mr{धास्ती}. Also an adverb.

\begin{cousinweb}
\textbf{\mr{बिनधास्त}} (\emph{bindhāsta}) is a Marathi compound of Hindustani \emph{bin}, \textquotedblleft{}without\textquotedblright{}, and \mr{धास्त}, \textquotedblleft{}dread\textquotedblright{}, the word in \mr{धास्ती}, whose origin is not settled.
\end{cousinweb}

\begin{grammarlens}[title={What you've built}]
One more word for describing what you lived and felt.
\end{grammarlens}

\subsection*{Your turn}
\begin{itemize}
\raggedright
  \item \emph{Say it:} \emph{bindhāsta}
  \item \emph{Say it:} \emph{bindhāsta}, once more
  \item \emph{Say it:} \emph{parīkṣā samplyāvar mulī bindhāsta kheḷlyā.} — When the exams were over the girls played without a care.
  \item \emph{Recall:} say the Marathi for adventurous, then the Marathi for chicken-hearted, then say \emph{bindhāsta} again
\end{itemize}

\begin{tcolorbox}[breakable,colback=teal!4,colframe=teal!35!black,title={Before you move on}]
What does \mr{बिनधास्त} mean? (\textquotedblleft{}Carefree; unafraid\textquotedblright{}.) Say it once more.
\end{tcolorbox}
\section[\texorpdfstring{\mr{बेफिकीर} (bephikīr)}{बेफिकीर (bephikīr)}]{\mr{बेफिकीर} (bephikīr) — unconcerned; heedless}

\label{lesson:MR-C475-bephikir}



\begin{quote}
Before the new one: say the Marathi for chicken-hearted, then the Marathi for carefree.
\end{quote}

\subsection*{You'll want to know: \mr{बेफिकीर}}
\textbf{\mr{बेफिकीर}} — \emph{bephikīr} — \textquotedblleft{}unconcerned; heedless\textquotedblright{}.

\begin{quote}
\mr{तो} \mr{खर्चाबद्दल} \mr{बेफिकीर} \mr{आहे}.
\end{quote}

\begin{quote}
\emph{to kharcābaddal bephikīr āhe.}
\end{quote}

\begin{quote}
He is heedless about spending.
\end{quote}

\mr{बे}-, without, + \mr{फिकीर}, worry. Unchanging.

\begin{cousinweb}
\textbf{\mr{बेफिकीर}} (\emph{bephikīr}) is a loan from Persian \emph{befikr}, \textquotedblleft{}without worry\textquotedblright{}, built on \emph{be-}, \textquotedblleft{}without\textquotedblright{}, and \emph{fikr}, \textquotedblleft{}thought, worry\textquotedblright{}, usually traced to Arabic.
\end{cousinweb}

\begin{grammarlens}[title={What you've built}]
One more word for describing what you lived and felt.
\end{grammarlens}

\subsection*{Your turn}
\begin{itemize}
\raggedright
  \item \emph{Say it:} \emph{bephikīr}
  \item \emph{Say it:} \emph{bephikīr}, once more
  \item \emph{Say it:} \emph{to kharcābaddal bephikīr āhe.} — He is heedless about spending.
  \item \emph{Recall:} say the Marathi for chicken-hearted, then the Marathi for carefree, then say \emph{bephikīr} again
\end{itemize}

\begin{tcolorbox}[breakable,colback=teal!4,colframe=teal!35!black,title={Before you move on}]
What does \mr{बेफिकीर} mean? (\textquotedblleft{}Unconcerned; heedless\textquotedblright{}.) Say it once more.
\end{tcolorbox}
\section[\texorpdfstring{\mr{भिणे} (bhiṇe)}{भिणे (bhiṇe)}]{\mr{भिणे} (bhiṇe) — to be scared of; to fear}

\label{lesson:MR-C475-bhine}



\begin{quote}
Before the new one: say the Marathi for carefree, then the Marathi for unconcerned.
\end{quote}

\subsection*{You'll want to know: \mr{भिणे}}
\textbf{\mr{भिणे}} — \emph{bhiṇe} — \textquotedblleft{}to be scared of; to fear\textquotedblright{}.

\begin{quote}
\mr{माझी} \mr{बहीण} \mr{कुत्र्यांना} \mr{भिते}.
\end{quote}

\begin{quote}
\emph{mājhī bahīṇ kutryānnā bhite.}
\end{quote}

\begin{quote}
My sister is afraid of dogs.
\end{quote}

Past (he …): \mr{तो} \mr{भ्याला}, \emph{to bhyālā}.

What is feared takes -\mr{ला}: \mr{कुत्र्याला} \mr{भिणे}. The past is irregular: \mr{तो} \mr{भ्याला}.

\begin{cousinweb}
\textbf{\mr{भिणे}} (\emph{bhiṇe}) is a tadbhava, inherited from Sanskrit \emph{bibheti}, \textquotedblleft{}fears\textquotedblright{}; CDIAL lists Marathi \emph{bhiṇẽ}, \textquotedblleft{}to fear\textquotedblright{} (CDIAL 9241).
\end{cousinweb}

\begin{grammarlens}[title={What you've built}]
One more word for describing what you lived and felt.
\end{grammarlens}

\subsection*{Your turn}
\begin{itemize}
\raggedright
  \item \emph{Say it:} \emph{bhiṇe}
  \item \emph{Say it:} \emph{bhiṇe}, once more
  \item \emph{Say it:} \emph{mājhī bahīṇ kutryānnā bhite.} — My sister is afraid of dogs.
  \item \emph{Recall:} say the Marathi for carefree, then the Marathi for unconcerned, then say \emph{bhiṇe} again
\end{itemize}

\begin{tcolorbox}[breakable,colback=teal!4,colframe=teal!35!black,title={Before you move on}]
What does \mr{भिणे} mean? (\textquotedblleft{}To be scared of; to fear\textquotedblright{}.) Say it once more.
\end{tcolorbox}
